\subsection{Changes to the Project}
\label{sec:alteracoes}

With respect to the initial plan, the representation changed from \num{60} descriptors with a dense network to graphs of the active channels, because the descriptors summarise the readout of the whole detector; the final models, however, still read \num{33} hand-crafted event-level scalars. Instead of a single architecture, the project compared \num{11} matched-capacity encoders, so that the number of parameters would not decide the result. Boosters fitted to aggregated features were added as reference models after a simple gradient-boosting model outperformed the GNN of the time in July 2026. The evaluation changed from a partition of \qty{80}{\percent} and \qty{20}{\percent}, with \num{10}-fold cross-validation, to a fixed stratified split of \num{7490}, \num{937} and \num{937} events, shared by all models; the final comparison had no cross-validation. APEX was compared at equal energy and also at equal detected light, because its simulation has an average light yield of \qty{180}{\pe\per\MeV}, \num{3.4} times that of the test set, \num{53.0}.

The extension to other energies and geometries, envisaged in the plan, was not carried out with the final models; what remained was the efficiency and energy sweep of the barycentre (Section~\ref{sec:limites}). The final comparison trained the encoders without the hyperparameter search of the plan (Section~\ref{sec:modelos}), which was applied only to the boosters (Section~\ref{sec:protocolo_comparacao}). Of the planned per-axis metrics, the report retains the RMSE, the mean of the three MAE values and $R^2$, the last in $z$ only (Section~\ref{sec:onde_erro}); the ranking and the training loss use the mean 3-D error (Section~\ref{sec:tarefa}). The comparison with APEX, proposed per coordinate, was restricted to $z$ and the combined width, because the correspondence of the transverse widths to $x$ and $y$ is not established (Section~\ref{sec:protocolo_comparacao}). Use in real-time data acquisition, expected in the plan, was not evaluated: the inference time was not measured, and the encoders, with capacity matched to \num{1843395} parameters, are larger than the network of the previous stage, with \num{8195}. That network is quoted rather than re-evaluated, because it used a different partition and the geometry prior to the corrections and because the trained model is not available.
